% Test that a note added between runs doesn't move the notes after it.
%
% The anchor page of a note is recorded in the .aux file and read in the
% next run.  If it were recorded under the note's number in the document, a
% note added in the second run (as when a citation becomes a footnote of
% its own) would give note B the number, and so the anchor page, of note C,
% which is called out on the next page, and B would be held back to that
% page.  Recorded under the callout's place in the source, B keeps its own.
%
% Three runs: the first records the anchor pages without note X, the second
% adds X, and the third checks, from the labels of the second, that note B
% was set on the page of its callout.
\documentclass{article}
\usepackage[a5paper,margin=18mm,marginparwidth=22mm]{geometry}
\usepackage{marginfix}
\usepackage{lipsum}

\makeatletter
\providecommand\akrun{0}
\AtBeginDocument{\if@filesw\immediate\write\@auxout{%
  \gdef\string\akrun{\the\numexpr\akrun+1\relax}}\fi}
% \akpage{<label>}: the page of a label, from the second field of \r@<label>
\def\akpage#1{%
  \expandafter\ifx\csname r@#1\endcsname\relax ??\else
    \expandafter\expandafter\expandafter\ak@page
      \csname r@#1\endcsname\@nil\fi}
\def\ak@page#1#2#3\@nil{#2}
\makeatother

\begin{document}
\ifnum\akrun>0
  Note X\marginpar{X: added in the second run.}
  is only there from the second run on.
\fi
\lipsum[1-2]

Note B\label{callout-B}\marginpar{B: its callout is on page 1.%
  \label{note-B}}
is called out on the first page.

\lipsum[3-6]

Note C\marginpar{C: its callout is on a later page.} is called out on a
later page.

\ifnum\akrun=2
  \edef\callout{\akpage{callout-B}}\edef\note{\akpage{note-B}}%
  \ifx\callout\note
    \typeout{anchorkeys test: note B on page \note, as its callout}%
  \else
    \PackageError{anchorkeys test}{Note B on page \note, its callout on
      page \callout}{}%
  \fi
\fi
\end{document}
